\section{Introduction}
\label{intro:section}

Virasoro constraints organize the descendant Gromov--Witten
invariants of a smooth projective variety into differential equations
for a single generating function. A basic structural question is
whether these equations pass from a base to the total space of a
geometric fibration. We prove that they pass through
projectivization of an arbitrary algebraic vector bundle.

For a smooth connected projective complex variety $Y$, let $Z_Y$
be its total ordinary descendant potential: the exponential of the
connected stable-map potentials, with genus weight $\hbar^{g-1}$,
ordinary cotangent-line classes at the markings, and monomials
$Q^d$ indexed by the actual effective classes
$d\in H_2(Y;\Z)$. We use all of $H^*(Y;\C)$ with its cohomological
parity and Poincar\'e pairing. The Virasoro operators are the
normally ordered operators associated to the first Hodge grading
\[
 \mu_Y|_{H^{p,q}(Y)}
       =\bigl(p-\tfrac12\dim_\C Y\bigr)\id,
 \qquad R_Y=c_1(TY)\cup.
\]
Their zero-mode constant is
$C_Y=\chi(Y)/16-\str(\mu_Y^2)/4$, and the other modes have no
scalar correction. Section~\ref{conv:section} specifies the
loop-space and quantization conventions completely. Write
$\V(Y)$ for the coefficientwise identities
\[
                      L_k^Y Z_Y=0\qquad(k\geq-1).
\]
Thus $\V(Y)$ includes every genus, individual integral curve
class, and finite list of descendants, including primitive and odd
insertions.

\begin{theorem}
\label{thm:main}
Let $B$ be a smooth connected projective variety over $\C$, and
let $E$ be an algebraic vector bundle of rank $r\geq2$ on $B$.
Let $X=\mathbb P_B(E)$ parametrize one-dimensional subspaces of
the fibers of $E$. Then
\[
                           \V(B)\ \Longrightarrow\ \V(X).
\]
\end{theorem}

The bundle is arbitrary: it need not split or admit a filtration
by line bundles, and it is subject to no positivity condition.
The base is allowed to have nonsemisimple quantum cohomology
and arbitrary Hodge types. In particular the assertion applies
successively to towers of projective bundles once the constraints
hold on the initial base. The conclusion concerns the ordinary
theory of each total space.

\subsection{Historical context and related results}

The Virasoro conjecture in Gromov--Witten theory extends the
differential constraints for intersection numbers on moduli
spaces of stable curves arising in Witten's conjecture and
Kontsevich's theorem \cite{Witten,Kontsevich}.
Eguchi--Hori--Xiong proposed the corresponding constraints
for quantum cohomology \cite{EHX}.
The extension to general Hodge types and odd cohomology,
including Katz's correction of the grading, is described by
Eguchi--Jinzenji--Xiong \cite{EJX}; see also
Getzler's account \cite{GetzlerVirasoro}.
Liu--Tian prove the constraints in genus zero \cite{LiuTian}.
In all genera, Givental's quantization formalism and Teleman's
reconstruction theorem establish them for targets with
generically semisimple quantum cohomology
\cite{GiventalQuant,GiventalSemisimple,Teleman}.
Okounkov--Pandharipande prove them for smooth target curves,
including odd insertions \cite{OkounkovPandharipande}.
Together with Theorem~\ref{thm:main}, this gives the full
constraints for every projective bundle over a smooth
projective curve, and for towers of such bundles.

The closest transfer theorem is due to
Coates--Givental--Tseng \cite{CGT}. They prove that the
Virasoro constraints hold for the base of a toric bundle if
and only if they hold for its total space, for toric bundles
constructed from sums of line bundles. Their proof combines
ancestor localization with Brown's mirror theorem
\cite{Brown}; split projective bundles are among its examples.
For arbitrary vector bundles, Fan \cite{Fan} proves all-genus
reconstruction and Chern-class dependence of projective-bundle
invariants. His argument uses the projective completion
$\mathbb P_B(E\oplus\mathcal O)$, which is also the auxiliary
space used below. Reconstruction of invariants by itself
does not establish compatibility with the Virasoro
differential operators.

Iritani--Koto \cite{IritaniKoto} construct a mirror theorem
and a decomposition of the quantum $D$-module for arbitrary
projective bundles. Koto \cite{Koto} proves a mirror theorem
for nonsplit toric bundles. These results concern genus zero.
In particular, Iritani--Koto's equivalence of generic
semisimplicity for a projective bundle and its base,
combined with Givental--Teleman, already gives the conclusion
of Theorem~\ref{thm:main} when the base has generically
semisimple quantum cohomology. The transfer proved here
also covers bases whose quantum cohomology is not
semisimple.

As in Fan's reconstruction, we use a projective completion
of the bundle as an auxiliary space.
The localization calculation adapts the ancestor organization
of Coates--Givental--Tseng, including their use of the
genus-zero cone and ancestor identities. We give the
additional gluing argument required when the orbit lines
are parametrized by a positive-dimensional variety.
The mechanism for passing constraints between the two fixed
components uses Iritani's equivariant shift operator
\cite{IritaniShift}. A local logarithm of that shift
cancels equivariant differentiation in the grading.
The resulting spectral modes can be quantized and their
equations continued between the components.
Quantum Riemann--Roch \cite{QRR} identifies the local
equations with the ordinary constraints. The argument
uses virtual localization \cite{GP} and the
ancestor--descendant formalism
\cite{KontsevichManin,Getzler,GiventalQuant};
the necessary forms of these results are recalled at
their points of use.

\subsection{The proof mechanism}

Twisting $E$ by a sufficiently negative line bundle leaves
$X$ unchanged and makes $E^*$ globally generated.
Set
\[
 W=\mathbb P_B(E\oplus\mathcal O).
\]
Let $\C^*$ scale the trivial summand, and let $\lambda$
be its equivariant parameter. Its fixed components are
$B$ and $X$. Write $y$ for the variable measuring
tautological degree and retain separate variables $Q^\beta$
for the full integral base classes.
The proof has five stages.

\begin{enumerate}
\item \emph{Express the auxiliary ancestors using the two
fixed theories.}
Localization factors the ancestor potential of $W$
through the product of the inverse-Euler-twisted ancestor
potentials of $B$ and $X$. The transformation is the
quantization of an upper symplectic series $R(z)$.
The same $R$ factors the genus-zero fundamental solution.
Section~\ref{loc:section} proves both statements together,
keeping the families of orbit lines as evaluation
correspondences. Localization provides a relation involving
both fixed theories; a second structure is needed to transfer
constraints from one to the other.

\item \emph{Separate the shift operator into spectral blocks.}
The genus-zero shift operator translates $\lambda$ by the
loop variable $z$. Its coefficients, at each base degree,
are polynomial in $y$. After reducing modulo $z$ and
positive base degree, its distinct spectral values are
$\lambda+h$, where
\[
                         h^r(h+\lambda)=y.
\]
There are $r+1$ branches at a generic value of $y$.
Near $y=0$, one tends to the eigenvalue zero and corresponds
to $B$; the other $r$ tend to $\lambda$ and together
correspond to $X$. Section~\ref{shift:section} establishes
this description without diagonalizing the quantum
cohomology of $B$.

\item \emph{Construct quantizable modes on the blocks.}
The equivariant grading contains
$\lambda\partial_\lambda$. On each block the logarithm
of the shift contains $z\partial_\lambda$.
Subtracting $\lambda/z$ times this logarithm removes
coefficient differentiation. The remaining grading defines
Virasoro-type loop operators.
Section~\ref{spec:section} constructs their projections
and logarithms coefficientwise and compares them at
$y=0$ with the fixed-component operators.
Section~\ref{fixed:section} uses quantum Riemann--Roch
to identify the latter, up to a triangular change of modes,
with the ordinary modes on $B$ or $X$.

\item \emph{Continue equations between branches.}
For fixed genus, base class, and insertions, the connected
ancestor invariants of $W$ are polynomial in $y$.
Moreover, ancestor powers are bounded independently of
fiber degree. Every coefficient of a quantized equation
therefore involves finitely many coefficients of the
spectral modes and is an actual identity of germs.
The covering $h\mapsto h^r(h+\lambda)$ is connected off
its branch values. Its monodromy transports the equations
from the $B$ branch to every branch. Adding the $r$
branches of the $X$ cluster and returning to $y=0$
gives the ordinary equations for $X$ up to scalar.

\item \emph{Fix the scalar normalization.}
Quantization is projective, so its covariance has only
determined the equations up to insertion-independent
scalars. The commutators
$[L_{-1},L_1]=-2L_0$ and $[L_0,L_k]=-kL_k$
remove those scalars and give exactly the prescribed
constant in $L_0$. Section~\ref{cont:section} proves
this calculation and completes the transfer.
\end{enumerate}

Two features of the argument are useful beyond the
localization formula itself. Subtracting the block logarithm
of a difference operator can turn a grading that
differentiates parameters into a loop operator over the
coefficient ring. Ancestor bounds then let identities for
such operators continue coefficientwise, even when
continuation of an entire quantized transformation has
not been defined. In the present setting these two steps
supply the passage from the ordinary constraints on one
fixed component to those on the other.
